\section{对后续队列的启发：为什么 EP139 应作为综述样板}

这期是 review/survey 类型，不是人物专访。它给后续Zhang Xiaojun AI 队列提供了一个处理模板：有 slides 时做 slide-complete；无 slides 时，不要反复贴人物帧，而要生成概念图、术语表和机制图。对综述类节目，笔记的价值在于重建知识结构，而不是压缩聊天内容。

\begin{knowledgebox}{综述类播客的生成模板}
第一，抽取技术谱系；第二，识别核心术语；第三，把口头判断转成表格或流程图；第四，标出争议和边界；第五，用总结章节连接到下一期或下一主题。这样得到的笔记才像讲义，而不是整理稿。
\end{knowledgebox}

\subsection{本章小结}

EP139 的方法论会用于后续广义 AI/互联网队列：只有封面和真正有信息的图进入正文；没有视觉内容时，用生成图表承载教学结构；每一期都保留 transcript、manifest、coverage 和 visual QA。


